\begin{resumo}[Abstract]
\begin{otherlanguage*}{english}
This document formalizes the Software Requirements Specification (SRS) for SIGEA-GTT (Intelligent Management System for Adverse Events -- Global Trigger Tool), a hospital computational platform dedicated to retrospective record review, clinical education, and adverse event surveillance, developed within the Department of Computing (DCOMP) and Department of Nursing at the Federal University of Sergipe (UFS). In strict compliance with ISO/IEC/IEEE 29148:2018 and the ISO/IEC 25010 product quality model, this specification rigorously details 25 Functional Requirements (RF01 to RF25) and 10 Non-Functional Requirements (RNF01 to RNF10). The functional requirements govern restricted institutional authentication, parameterization of the Institute for Healthcare Improvement (IHI) clinical trigger catalog, NCC MERP harm severity scoring, academic class management, high-fidelity clinical record simulation in JSONB format, retrospective auditing under a strict 20-minute countdown timer, seamless integration of six Quality Management Tools (Ishikawa 6M diagram, dynamic GUT prioritization matrix, 5W2H action matrix, PDCA cycle, SWOT matrix, and ideation board), and automated calculation of hospital epidemiological and pedagogical performance indicators. The non-functional requirements enforce strict desktop clinical ergonomics across homologated screen resolutions (WUXGA, HD+, HD, and 21:9 Ultrawide), REST API response latency under 200 milliseconds, cryptographic security via stateless JWT and BCrypt hashing, full clinical data isolation in compliance with the Brazilian General Data Protection Law (LGPD), and unyielding Quality Gates mandating 100\% automated test coverage and 100\% codebase documentation. Additionally, the document provides the UML 2.5 use case diagram and a bidirectional requirements traceability matrix.

\vspace{\onelineskip}
\noindent
\textbf{Keywords}: Requirements Engineering. Software Requirements Specification. Global Trigger Tool. Patient Safety. Health Information Systems. Federal University of Sergipe.
\end{otherlanguage*}
\end{resumo}
